\chapter{Presupuesto}
\label{cap:presupuesto}

Este cap\'itulo cuantifica el coste de ejecuci\'on del proyecto. El presupuesto recoge
cuatro partidas: material fungible y suministros, amortizaci\'on de los equipos, licencias
de software y mano de obra. A ellas se a\~naden unos costes indirectos y, con todo, se
obtiene el importe total.

Tres criterios de imputaci\'on gobiernan las cifras que siguen. En primer lugar, el
periodo imputado abarca siete meses, de marzo a septiembre de 2026, que comprenden
desde la revisi\'on bibliogr\'afica hasta la entrega de la memoria. En segundo lugar, los
equipos se amortizan de forma lineal y se imputa la fracci\'on correspondiente a esos
siete meses, con independencia de su uso efectivo dentro de cada jornada. En tercer lugar,
los importes se expresan sin impuestos indirectos, puesto que el trabajo se desarrolla en
el \'ambito acad\'emico y no da lugar a facturaci\'on.

\section{Material fungible y suministros}
\label{sec:fungible}

El proyecto es de naturaleza computacional, de modo que apenas consume material f\'isico.
La \'unica partida relevante es la energ\'ia el\'ectrica, dominada por los entrenamientos
desatendidos. Seg\'un el \cref{sec:coste}, estos ocuparon \SI{237.1}{\hour} de c\'omputo;
con un consumo aproximado de \SI{150}{\watt} bajo carga, suman unos \num{36}
kilovatios-hora. El resto de la dedicaci\'on a\~nade otros \num{24} kilovatios-hora, con el
equipo en r\'egimen de trabajo ordinario. El precio de \num{0.15}~\texteuro{} por
kilovatio-hora es una tarifa convencional de imputaci\'on adoptada para el proyecto, no un
precio contratado. La \cref{tab:fungible} recoge ambos conceptos.

\begin{table}[htbp]
  \centering
  \caption{Presupuesto de material fungible y suministros.}
  \label{tab:fungible}
  \small
  \begin{tabular}{lrrr}
    \toprule
    Concepto & Cantidad & Coste unitario (\texteuro) & Total (\texteuro) \\
    \midrule
    Energ\'ia el\'ectrica (kWh)        & \num{60} & \num{0.15} & \num{9.00} \\
    Material de oficina e impresi\'on  & \num{1}  & \num{45.00} & \num{45.00} \\
    \midrule
    \textbf{Total} & & & \textbf{\num{54.00}} \\
    \bottomrule
  \end{tabular}
  \notatabla{Se adopta un precio de \num{0.15} \texteuro/kWh, valor de referencia para el
    suministro dom\'estico durante el periodo imputado.}
  \fuente{elaboraci\'on propia}
\end{table}

\section{Amortizaci\'on de los equipos}
\label{sec:equipos}

El proyecto se ejecuta sobre un \'unico equipo, descrito en el \cref{sec:entorno}. Dado
que su vida \'util excede la duraci\'on del trabajo, no se imputa el precio de
adquisici\'on sino la amortizaci\'on correspondiente al periodo de uso, seg\'un la
\cref{eq:amortizacion}.

\begin{equation}
  A = C \cdot \frac{t}{V}
  \label{eq:amortizacion}
\end{equation}

En la \cref{eq:amortizacion}, $A$ es el importe imputado, $C$ el precio de adquisici\'on,
$t$ el periodo de uso y $V$ la vida \'util, ambos expresados en meses. Se aplica
amortizaci\'on lineal y sin valor residual, criterio habitual para equipos de proceso de
informaci\'on. La \cref{tab:equipos} detalla el resultado para cada elemento.

\begin{table}[htbp]
  \centering
  \caption{Amortizaci\'on de los equipos imputada al proyecto.}
  \label{tab:equipos}
  \small
  \begin{tabular}{lrrrr}
    \toprule
    Equipo & Precio (\texteuro) & Vida \'util (a\~nos) & Uso (meses) &
      Imputaci\'on (\texteuro) \\
    \midrule
    Port\'atil con GPU dedicada & \num{1800.00} & \num{4} & \num{7} & \num{262.50} \\
    Monitor externo                & \num{180.00}  & \num{5} & \num{7} & \num{21.00} \\
    Unidad externa de respaldo     & \num{60.00}   & \num{3} & \num{7} & \num{11.67} \\
    \midrule
    \textbf{Total} & & & & \textbf{\num{295.17}} \\
    \bottomrule
  \end{tabular}
  \fuente{elaboraci\'on propia}
\end{table}

Conviene se\~nalar que la unidad de procesamiento gr\'afico condiciona el dise\~no
experimental, seg\'un se argument\'o en el \cref{sec:entorno}, pero apenas condiciona el
presupuesto. Su amortizaci\'on forma parte de los \num{262.50} \texteuro{} imputados al
port\'atil, importe menor que el coste de dos jornadas de trabajo.

\section{Licencias de software}
\label{sec:software}

Todas las herramientas empleadas son de c\'odigo abierto o de uso gratuito, de modo que
esta partida asciende a \num{0.00} \texteuro. La \cref{tab:componentes} enumera la pila de
software: Python, PyTorch, torchvision, timm, diffusers, robomimic, Hydra y Zarr se
distribuyen bajo licencias libres. A ellas se suman la distribuci\'on Ubuntu, el sistema
de composici\'on \TeX{} Live y el control de versiones Git, tambi\'en libres. La licencia de
Windows 11 viene incluida en el precio del port\'atil y queda, por tanto, recogida en el
\cref{sec:equipos}.

\section{Mano de obra}
\label{sec:mano-de-obra}

La mano de obra distingue dos perfiles. El autor asume la implementaci\'on, la
ejecuci\'on, el an\'alisis y la redacci\'on, con el coste horario de un ingeniero de datos
junior. El director asume la definici\'on del alcance, el seguimiento peri\'odico y la
revisi\'on de la memoria, con el coste horario de un investigador s\'enior. Ambos importes
corresponden al coste bruto para la entidad, incluidas las cargas sociales. Conviene declarar
su origen: son tarifas convencionales de imputaci\'on adoptadas para el proyecto, de un orden
habitual para esos dos perfiles, y no proceden de una tabla salarial ni de una fuente
retributiva concreta. El
presupuesto es, en consecuencia, una estimaci\'on del orden de magnitud del coste y no una
oferta; la mano de obra domina el total, de modo que cualquier revisi\'on de esas tarifas se
traslada casi \'integramente a la cifra final.

La \cref{tab:horas} desglosa la dedicaci\'on del autor por fases. El reparto refleja el
peso real de cada actividad: la preparaci\'on del entorno de c\'omputo result\'o m\'as
costosa de lo previsto, mientras que la ejecuci\'on de los entrenamientos, aunque ocup\'o
m\'as de \num{237} horas de m\'aquina, exigi\'o poca intervenci\'on por tratarse de procesos
desatendidos.

\begin{table}[htbp]
  \centering
  \caption{Dedicaci\'on del autor por fases del proyecto.}
  \label{tab:horas}
  \small
  \begin{tabular}{lr}
    \toprule
    Fase & Horas \\
    \midrule
    Revisi\'on bibliogr\'afica y estado del arte          & \num{70} \\
    Preparaci\'on del entorno de c\'omputo                & \num{55} \\
    Implementaci\'on de las variantes del codificador     & \num{90} \\
    Ejecuci\'on y supervisi\'on de los entrenamientos     & \num{60} \\
    An\'alisis de resultados y comprobaci\'on cualitativa & \num{65} \\
    Redacci\'on de la memoria                             & \num{110} \\
    \midrule
    \textbf{Total} & \textbf{\num{450}} \\
    \bottomrule
  \end{tabular}
  \fuente{elaboraci\'on propia}
\end{table}

La dedicaci\'on del director se estima en \num{40} horas: unas \num{20} reuniones de
seguimiento de una hora, repartidas a lo largo de los siete meses, y otras \num{20} horas
de revisi\'on de los borradores. La \cref{tab:mano-de-obra} aplica los costes horarios a
ambos perfiles.

\begin{table}[htbp]
  \centering
  \caption{Presupuesto de mano de obra.}
  \label{tab:mano-de-obra}
  \small
  \begin{tabular}{lrrr}
    \toprule
    Perfil & Horas & Coste horario (\texteuro) & Total (\texteuro) \\
    \midrule
    Ingeniero de datos junior (autor) & \num{450} & \num{30.00} & \num{13500.00} \\
    Investigador s\'enior (director)  & \num{40}  & \num{60.00} & \num{2400.00} \\
    \midrule
    \textbf{Total} & \textbf{\num{490}} & & \textbf{\num{15900.00}} \\
    \bottomrule
  \end{tabular}
  \fuente{elaboraci\'on propia}
\end{table}

\section{Resumen del presupuesto}
\label{sec:resumen-presupuesto}

La \cref{tab:resumen-presupuesto} agrega las cuatro partidas anteriores y a\~nade unos
costes indirectos del \SI{15}{\percent} sobre el subtotal, porcentaje que cubre espacio de
trabajo, conectividad y servicios generales no imputables a un concepto concreto.

\begin{table}[htbp]
  \centering
  \caption{Resumen del presupuesto.}
  \label{tab:resumen-presupuesto}
  \small
  \begin{tabular}{lr}
    \toprule
    Partida & Importe (\texteuro) \\
    \midrule
    Material fungible y suministros & \num{54.00} \\
    Amortizaci\'on de equipos       & \num{295.17} \\
    Licencias de software           & \num{0.00} \\
    Mano de obra                    & \num{15900.00} \\
    \midrule
    Subtotal                        & \num{16249.17} \\
    Costes indirectos (\SI{15}{\percent}) & \num{2437.38} \\
    \midrule
    \textbf{Total}                  & \textbf{\num{18686.55}} \\
    \bottomrule
  \end{tabular}
  \fuente{elaboraci\'on propia}
\end{table}

El coste total del proyecto asciende, por tanto, a \num{18686.55} \texteuro. La
distribuci\'on entre partidas resulta muy desigual: la mano de obra representa el
\SI{85}{\percent} del total, mientras que los equipos apenas alcanzan el
\SI{1.6}{\percent}. Esa asimetr\'ia es caracter\'istica de un proyecto de investigaci\'on
computacional ejecutado sobre hardware propio.

De ah\'i se sigue una consideraci\'on econ\'omica sobre el dise\~no experimental. Las
limitaciones descritas en el \cref{sec:coste} no proceden del presupuesto, sino del tiempo
de calendario: ampliar el estudio a tres semillas por variante multiplicar\'ia por tres el
c\'omputo sin apenas alterar el coste, ya que la amortizaci\'on y la energ\'ia son
partidas menores. El factor limitante es que esas horas de m\'aquina deben transcurrir
dentro del plazo del proyecto, no lo que cuestan.
